\documentclass[11pt]{article}
\usepackage[margin=.8in]{geometry}\usepackage[T1]{fontenc}\usepackage{lmodern,microtype}
\usepackage{amsmath,booktabs,graphicx,pdflscape}\usepackage[hidelinks]{hyperref}
\setlength{\parindent}{0pt}\setlength{\parskip}{6pt}\setlength{\emergencystretch}{2em}
\title{Four-year joint refit and fixed-parameter sensitivity\\Reference C-to-D probability $\mu=0.01$}
\author{Model, completed results and patient analyses}\date{5 November 2026}
\begin{document}\maketitle
\section{Reference model and unchanged cohort}
The reference now fixes $\mu=0.01$ per C division and $D_{\rm init}=0$ at the measured baseline. Mutation applies before and after CAR. All pre-CAR fitnesses are 1. After CAR, $f_W=1$, $f_C=1+g(t;\alpha,\theta)$, and $f_D=f_C+s_D$. The Gamma density integrates to 1 on $[0,\infty)$ for every shape and scale; its peak is derived.
\[
g(t;\alpha,\theta)=\frac{t^{\alpha-1}e^{-t/\theta}}{\Gamma(\alpha)\theta^\alpha},\qquad t\geq0,\qquad
\int_0^\infty g(t;\alpha,\theta)\,dt=1.
\]
Alpha is the shape and theta is the scale in days. Time is supplied numerically in days, and fitness uses the numerical density value with a fixed multiplier of 1. The legacy code field \texttt{area\_days} names that multiplier; it does not mean that the normalized density integral is one day. Mu is an effective second-hit probability, not a nucleotide mutation rate.
Fix $N=100000$, $b=1/28$ per day, $s_D=0.02$, the strict threshold $D/N>0.01$ (1001 cells), zero diagnosis delay, and observation SD 0.25 on the $\operatorname{logit}(2VAF)$ scale.
Refit one cohort-wide $\alpha\in[2,50]$ and $\theta\in[1,1460]$ days with the same shared-state VAF/tMN joint likelihood.
The frozen cohort has 14 patients, 42 follow-up VAF observations, exact events at days 127 and 589, one censor record at day 1460, and 11 unknown outcomes. Unknown outcomes are not negatives.
\section{Joint likelihood}
The founder marker is retained after conversion from C to D, so the latent VAF is $(C+D)/(2N)$. Follow-up measurements have Gaussian error on the $\operatorname{logit}(2VAF)$ scale. The shared-state forward recursion weights the D distribution at every VAF measurement, then evaluates first-passage boundary flux at an exact tMN time or surviving probability at a censoring time. Patients with unknown clinical outcomes contribute VAF likelihood only. Baselines initialize the process and are not counted again as follow-up observations. The reported VAF and clinical-given-VAF terms decompose this joint likelihood; they are not independent marginal fits.
\section{Initial states}
Preserve total founder burden $C_i+D_i=2N\,VAF_{i,0}$ at the measured baseline. Subtract the chosen initial D count from that total to initialize C; W makes up N. Use the original continuous initial fraction in the hybrid inference and integer rounding for full stochastic simulations.
The smallest baseline founder has approximately 544 cells. The common D grid is therefore capped at 500 cells. Evolve negative baseline dates neutrally with mutation to event-free CAR day 0.
\section{Sensitivity design}
\begin{center}\begin{tabular}{lp{.75\linewidth}}\toprule Parameter & Levels\\\midrule
$\mu$ & $0,10^{-6},10^{-5},10^{-4},0.001,0.003,0.01,0.03,0.1$\\
$D_{\rm init}$ & $0,1,3,10,30,100,300,500$ cells\\
$b$ & $1/280,1/84,1/28,1/7$ per day\\
$s_D$ & $0,0.01,0.02,0.05,0.1$\\\bottomrule\end{tabular}\end{center}
The complete mu-by-D grid contains 72 scenarios. The zero-mutation/zero-D corner is structurally incompatible with positive tMN events and is recorded without a probability floor. The b and s sensitivities vary one parameter at a time at reference mu and D, adding seven scenarios. Thus 79 scenarios imply 78 compatible joint fits. All compatible scenarios re-estimate alpha and theta; neither is held at the reference estimate.
The levels are model-sensitivity choices, not biological confidence intervals. Slow turnover endpoints have precedent in \href{https://pmc.ncbi.nlm.nih.gov/articles/PMC8109274/}{Heyde et al., Cell (2021)}; $1/7$ is a fast-turnover stress case.
\section{Numerical validation before batch fitting}
At higher mu the D distribution need not be rare. Preserve the existing deterministic-q/stochastic-D inference model and use a conservative moment-based proposal for each unknown-outcome numerical state cap. The proposal only chooses the finite-N chain's reflecting truncation; the actual likelihood still evolves its state probabilities. Compare larger caps and solver tolerances on log likelihood; check the reference against the full D=0,...,N chain.
Exponential tilting is an invertible numerical change of variables. Its initial scale must account for immigration as well as initial D to prevent high-mu moment overflow. Rescale the killed chain more frequently in high-mu stress cases. Verify event densities independently by sparse matrix exponentiation for low and high immigration and seeded D.
Use multiple optimizer starts and retain all evaluations, bounds, convergence status, numerical checks and source hashes. Reuse a checkpoint only when parameters, cohort, bounds and inference source agree.
\section{Predictions, figures and delivery}
For each fitted scenario calculate each patient's year 1--4 first-passage probabilities from the own baseline, conditional on event-free CAR day 0. Subsequent individual VAF and eventual outcome do not set baseline-initialized forecasts.
Produce joint likelihood, fitted alpha/theta and patient risk heatmaps over mu and D, together with b and s comparisons. Update the new-reference patient VAF, tMN and trajectory figures, using 200 full Moran paths per patient (2800 total). Compare full Moran process means and event proportions with the inference approximation.
Write code first, then this English TeX protocol, then compute. Preserve the previous study and export English PDF/TeX, XLSX/CSV, code, seeds and raw reference trajectories. All predictions are conditional on fixed scenarios; the aggregate threshold is not a clinically calibrated diagnosis model.
\IfFileExists{executed_results.tex}{\clearpage\section{Executed joint reference fit}

The reference fixes mu=0.01 and initial D=0. Its fitted shape is 15.242771, scale 1 days, and derived mode 14.242771 days. Joint NLL is 455.808076, decomposed into VAF NLL 291.949583 and clinical-given-VAF NLL 163.858493. The preceding mu=1e-6 reference had NLL 511.718830. Parameters were re-estimated after changing mu.

The lowest NLL across compatible sensitivity scenarios is 308.372238, in mu\_0\_D\_500. This is a conditional scenario comparison; the requested reference remains mu=0.01 and initial D=0.

All 78 conditional optimizations ended at at least one search boundary. At the reference, the 95\% one-step observation interval covers 22/42 follow-up VAFs (52.4\%). VAF RMSE is 4.856 percentage points. These diagnostics show substantial model/data mismatch; the fitted timescale and clinical probabilities should be interpreted conditionally, rather than as validated physiological estimates.

\begin{landscape}\section{Mu and initial-D sensitivity}\begin{center}\includegraphics[width=.98\linewidth,height=.78\textheight,keepaspectratio]{figures/mu_D_joint_likelihood.pdf}\end{center}

Every compatible cell re-optimizes alpha and theta. The gray zero-mutation/zero-D corner is structurally incompatible with observed positive tMN. No probability floor is used.

\clearpage\begin{center}\includegraphics[width=\linewidth]{figures/mu_D_fitted_parameters.pdf}\end{center}

Parameter maps show conditional optima within the common bounds. A value at a search boundary is not an identified physiological timescale. At alpha=50 and theta=1460 days, the Gamma mode is 71540 days. The fraction of the normalized Gamma area occurring within the four-year observation window therefore varies across scenarios, while its full integral remains 1.

\end{landscape}\clearpage\section{Turnover and double-hit fitness sensitivity}

\begin{center}\includegraphics[width=\linewidth]{figures/b_sensitivity.pdf}\\[10pt]\includegraphics[width=\linewidth]{figures/s_D_sensitivity.pdf}\end{center}

\clearpage\section{New-reference patient diagnostics}

\begin{center}\includegraphics[width=.98\linewidth]{patient_analysis/figures/fit_diagnostics.pdf}\end{center}

\clearpage\section{New-reference Gamma pulse and one-year risk}

\begin{center}\includegraphics[width=\linewidth]{patient_analysis/figures/gamma_and_fitness.pdf}\\[10pt]\includegraphics[width=.88\linewidth,height=.65\textheight,keepaspectratio]{patient_analysis/figures/one_year_tMN_heatmap.pdf}\end{center}

\clearpage\subsection{Reference patient risk table}

\begin{center}\small\begin{tabular}{lrrrr}\toprule Patient & Year 1 & Year 2 & Year 3 & Year 4\\\midrule

Patient\_25 & $2.349\times10^{-5}$ & $7.595\times10^{-1}$ & $9.989\times10^{-1}$ & $10.000\times10^{-1}$\\

Patient\_44 & $8.728\times10^{-5}$ & $8.375\times10^{-1}$ & $9.996\times10^{-1}$ & $10.000\times10^{-1}$\\

Patient\_102 & $1.056\times10^{-3}$ & $9.649\times10^{-1}$ & $10.000\times10^{-1}$ & $10.000\times10^{-1}$\\

Patient\_103 & $2.525\times10^{-20}$ & $6.532\times10^{-8}$ & $4.381\times10^{-4}$ & $1.918\times10^{-2}$\\

Patient\_123 & $4.014\times10^{-1}$ & $10.000\times10^{-1}$ & $10.000\times10^{-1}$ & $10.000\times10^{-1}$\\

Patient\_180 & $3.240\times10^{-14}$ & $4.256\times10^{-4}$ & $1.071\times10^{-1}$ & $5.416\times10^{-1}$\\

Patient\_197 & $1.093\times10^{-11}$ & $1.253\times10^{-2}$ & $4.865\times10^{-1}$ & $9.153\times10^{-1}$\\

Patient\_235 & $10.000\times10^{-1}$ & $10.000\times10^{-1}$ & $10.000\times10^{-1}$ & $10.000\times10^{-1}$\\

Patient\_238 & $5.991\times10^{-2}$ & $9.998\times10^{-1}$ & $10.000\times10^{-1}$ & $10.000\times10^{-1}$\\

Patient\_316 & $10.000\times10^{-1}$ & $10.000\times10^{-1}$ & $10.000\times10^{-1}$ & $10.000\times10^{-1}$\\

Patient\_318 & $10.000\times10^{-1}$ & $10.000\times10^{-1}$ & $10.000\times10^{-1}$ & $10.000\times10^{-1}$\\

Patient\_321 & $9.410\times10^{-1}$ & $10.000\times10^{-1}$ & $10.000\times10^{-1}$ & $10.000\times10^{-1}$\\

Patient\_356 & $2.191\times10^{-1}$ & $10.000\times10^{-1}$ & $10.000\times10^{-1}$ & $10.000\times10^{-1}$\\

Patient\_445 & $1.685\times10^{-21}$ & $6.084\times10^{-8}$ & $5.606\times10^{-4}$ & $2.487\times10^{-2}$\\

\bottomrule\end{tabular}\end{center}

Risk forecasts condition on each patient's measured baseline and event-free CAR day 0. They do not use the subsequent individual VAF or eventual outcome to set the forecast. The fitted common parameters were learned in this cohort, so these are in-sample conditional predictions.

\subsection{Full-process simulation comparison}

Across 56 patient/year comparisons at the new reference, the maximum difference between the inference approximation and 200 full Moran paths is 0.0982 percentage points for mean C and 0.0758 percentage points for mean D. The largest absolute tMN probability difference is 0.0336. All 56 model probabilities fall within pointwise exact 95\% binomial intervals for the simulated event proportions. These intervals measure Monte Carlo error at fixed parameters, not fitted-parameter uncertainty. With 200 paths, empirical event fractions have increments of 0.005 and cannot verify extremely small tail probabilities. Agreement with the simulator verifies the approximation in this reference setting; it does not establish agreement with clinical incidence.

\subsection{Patient figure guide}

Open diamonds are baseline VAF; black points are follow-up measurements. The blue VAF band is a 95\% observation prediction interval. Blue C and red D trajectory panels contain all 200 full Moran paths with opacity 0.035 and their arithmetic mean. Dotted lines show known diagnosis or censoring times. Process paths continue after diagnosis without modeling postdiagnosis treatment. Unknown outcomes remain unknown.

\clearpage\begin{center}\includegraphics[width=.97\linewidth,height=.9\textheight,keepaspectratio]{patient_analysis/figures/patients/Patient_25_dashboard.pdf}\end{center}

\clearpage\begin{center}\includegraphics[width=.97\linewidth,height=.9\textheight,keepaspectratio]{patient_analysis/figures/patients/Patient_44_dashboard.pdf}\end{center}

\clearpage\begin{center}\includegraphics[width=.97\linewidth,height=.9\textheight,keepaspectratio]{patient_analysis/figures/patients/Patient_102_dashboard.pdf}\end{center}

\clearpage\begin{center}\includegraphics[width=.97\linewidth,height=.9\textheight,keepaspectratio]{patient_analysis/figures/patients/Patient_103_dashboard.pdf}\end{center}

\clearpage\begin{center}\includegraphics[width=.97\linewidth,height=.9\textheight,keepaspectratio]{patient_analysis/figures/patients/Patient_123_dashboard.pdf}\end{center}

\clearpage\begin{center}\includegraphics[width=.97\linewidth,height=.9\textheight,keepaspectratio]{patient_analysis/figures/patients/Patient_180_dashboard.pdf}\end{center}

\clearpage\begin{center}\includegraphics[width=.97\linewidth,height=.9\textheight,keepaspectratio]{patient_analysis/figures/patients/Patient_197_dashboard.pdf}\end{center}

\clearpage\begin{center}\includegraphics[width=.97\linewidth,height=.9\textheight,keepaspectratio]{patient_analysis/figures/patients/Patient_235_dashboard.pdf}\end{center}

\clearpage\begin{center}\includegraphics[width=.97\linewidth,height=.9\textheight,keepaspectratio]{patient_analysis/figures/patients/Patient_238_dashboard.pdf}\end{center}

\clearpage\begin{center}\includegraphics[width=.97\linewidth,height=.9\textheight,keepaspectratio]{patient_analysis/figures/patients/Patient_316_dashboard.pdf}\end{center}

\clearpage\begin{center}\includegraphics[width=.97\linewidth,height=.9\textheight,keepaspectratio]{patient_analysis/figures/patients/Patient_318_dashboard.pdf}\end{center}

\clearpage\begin{center}\includegraphics[width=.97\linewidth,height=.9\textheight,keepaspectratio]{patient_analysis/figures/patients/Patient_321_dashboard.pdf}\end{center}

\clearpage\begin{center}\includegraphics[width=.97\linewidth,height=.9\textheight,keepaspectratio]{patient_analysis/figures/patients/Patient_356_dashboard.pdf}\end{center}

\clearpage\begin{center}\includegraphics[width=.97\linewidth,height=.9\textheight,keepaspectratio]{patient_analysis/figures/patients/Patient_445_dashboard.pdf}\end{center}

\clearpage\section{Patient-specific mu--D risk responses}

Each patient's sensitivity panel holds the measured baseline founder burden fixed and subtracts D from that total when initializing C. Colors are probabilities on logarithmic scales. Year-one and year-four maps are supplied for every patient.

\begin{center}\includegraphics[width=.98\linewidth,height=.85\textheight,keepaspectratio]{figures/patients/Patient_25_mu_D_risk.pdf}\end{center}

\clearpage\begin{center}\includegraphics[width=.98\linewidth,height=.85\textheight,keepaspectratio]{figures/patients/Patient_44_mu_D_risk.pdf}\end{center}

\clearpage\begin{center}\includegraphics[width=.98\linewidth,height=.85\textheight,keepaspectratio]{figures/patients/Patient_102_mu_D_risk.pdf}\end{center}

\clearpage\begin{center}\includegraphics[width=.98\linewidth,height=.85\textheight,keepaspectratio]{figures/patients/Patient_103_mu_D_risk.pdf}\end{center}

\clearpage\begin{center}\includegraphics[width=.98\linewidth,height=.85\textheight,keepaspectratio]{figures/patients/Patient_123_mu_D_risk.pdf}\end{center}

\clearpage\begin{center}\includegraphics[width=.98\linewidth,height=.85\textheight,keepaspectratio]{figures/patients/Patient_180_mu_D_risk.pdf}\end{center}

\clearpage\begin{center}\includegraphics[width=.98\linewidth,height=.85\textheight,keepaspectratio]{figures/patients/Patient_197_mu_D_risk.pdf}\end{center}

\clearpage\begin{center}\includegraphics[width=.98\linewidth,height=.85\textheight,keepaspectratio]{figures/patients/Patient_235_mu_D_risk.pdf}\end{center}

\clearpage\begin{center}\includegraphics[width=.98\linewidth,height=.85\textheight,keepaspectratio]{figures/patients/Patient_238_mu_D_risk.pdf}\end{center}

\clearpage\begin{center}\includegraphics[width=.98\linewidth,height=.85\textheight,keepaspectratio]{figures/patients/Patient_316_mu_D_risk.pdf}\end{center}

\clearpage\begin{center}\includegraphics[width=.98\linewidth,height=.85\textheight,keepaspectratio]{figures/patients/Patient_318_mu_D_risk.pdf}\end{center}

\clearpage\begin{center}\includegraphics[width=.98\linewidth,height=.85\textheight,keepaspectratio]{figures/patients/Patient_321_mu_D_risk.pdf}\end{center}

\clearpage\begin{center}\includegraphics[width=.98\linewidth,height=.85\textheight,keepaspectratio]{figures/patients/Patient_356_mu_D_risk.pdf}\end{center}

\clearpage\begin{center}\includegraphics[width=.98\linewidth,height=.85\textheight,keepaspectratio]{figures/patients/Patient_445_mu_D_risk.pdf}\end{center}
}{}
\end{document}
